\documentclass[a4paper,12pt]{article}

% --------------------------------------------------------
% CODIFICAÇÃO E TIPOGRAFIA
% --------------------------------------------------------
\usepackage[utf8]{inputenc}       % Permite acentuação direta no código-fonte
\usepackage[T1]{fontenc}          % Usa codificação T1 (melhor para português)
\usepackage[brazil]{babel}        % Tradução e hifenização para português
\usepackage{microtype}            % Melhora o espaçamento entre letras/palavras
\usepackage{parskip}              % Espaço entre parágrafos em vez de recuo

% --------------------------------------------------------
% PACOTES MATEMÁTICOS
% --------------------------------------------------------
\usepackage{amsmath}              % Ambientes matemáticos como align, equation
%\usepackage{amssymb}              % Símbolos matemáticos adicionais
\usepackage{amsthm}               % Ambientes de teoremas
\usepackage{mathtools}            % Extensões do amsmath (ex: \coloneqq)
\usepackage{mathrsfs}             % Fonte caligráfica com \mathscr
\usepackage{bbm}                  % Indicador \mathbbm{1}
\usepackage{dsfont}               % Alternativa para conjuntos: \mathds
\usepackage{commath}              % Notações como \abs{}, \norm{}, \eval{}
\usepackage{pgfmath}             % Permite realizar cálculos matemáticos (somas, produtos, números aleatórios, etc.)
\usepackage{siunitx}              % sistema de unidades
\sisetup{group-separator = {\,}} % espaço fino

% --------------------------------------------------------
% FONTE MODERNA (TEXTO E MATEMÁTICA)
% --------------------------------------------------------
%\usepackage{newtxtext}            % Fonte do texto (estilo Times)
%\usepackage{newtxmath}            % Fonte matemática compatível com pacotes AMS
\usepackage[bitstream-charter]{mathdesign} % Fonte elegante e matemática harmoniosa

% --------------------------------------------------------
% AMBIENTES TEÓRICOS PERSONALIZADOS
% --------------------------------------------------------
\theoremstyle{definition}
\newtheorem{definicao}{Definição}       % Definição numerada por seção
\newtheorem{exemplo}{Exemplo}           % Exemplo numerado por seção

%\theoremstyle{plain}
\newtheorem{teorema}{Teorema}           % Teorema numerado por seção
\newtheorem{lema}[teorema]{Lema}                 % Lema com numeração conjunta
\newtheorem{proposicao}[teorema]{Proposição}     % Proposição idem
\newtheorem{corolario}[teorema]{Corolário}       % Corolário idem
\newtheorem{exercicio}[teorema]{Exercício}

%\theoremstyle{remark}
\newtheorem{observacao}[teorema]{Observação}     % Observação com mesmo contador

% --------------------------------------------------------
% TABELAS E FIGURAS
% --------------------------------------------------------
\usepackage{graphicx}             % Inclusão de imagens
\usepackage{float}                % Controle de posicionamento (ex: [H])
\usepackage{caption}              % Personalização de legendas
\usepackage{subcaption}          % Subfiguras com \begin{subfigure}
\usepackage{booktabs}            % Tabelas com qualidade profissional
\usepackage{array}               % Mais opções em colunas de tabelas
\usepackage{multirow}            % Células que ocupam várias linhas
\usepackage{multicol}            % Disposição em colunas múltiplas
\usepackage{tabularx}            % Tabelas com colunas de largura ajustável
\usepackage{longtable}           % Tabelas que quebram página
\usepackage{arydshln}            % Linhas tracejadas horizontais e verticais
\usepackage{makecell}            % Células com múltiplas linhas (e quebra de linha)
\usepackage{diagbox}             % Cabeçalho diagonal em tabelas
\usepackage{pdflscape}           % Gira páginas (modo paisagem)
\newcolumntype{L}[1]{>{\raggedright\let\newline\\\arraybackslash\hspace{0pt}}m{#1}}
\newcolumntype{C}[1]{>{\centering\let\newline\\\arraybackslash\hspace{0pt}}m{#1}}
\newcolumntype{R}[1]{>{\raggedleft\let\newline\\\arraybackslash\hspace{0pt}}m{#1}}

% --------------------------------------------------------
% FORMATAÇÃO E ESTRUTURA
% --------------------------------------------------------
\usepackage{geometry}            % Configuração de margens
\geometry{a4paper, left=1cm, right=1cm, bottom=1.5cm, top=1.5cm, headsep=1cm, footskip=1cm}
\usepackage{titlesec}            % Controle sobre títulos de seções
\usepackage{fancyhdr}            % Cabeçalhos e rodapés personalizados
\pagestyle{fancy}
\fancyhf{}                % Limpa cabeçalho e rodapé
\fancyfoot[R]{\thepage}   % Numeração à direita no rodapé
\renewcommand{\headrulewidth}{0pt} % Remove linha no topo
\usepackage[shortlabels]{enumitem}            % Listas com controle de espaçamento e símbolo
\usepackage{tikz}                % Desenho de gráficos vetoriais (diagramas, etc.)
\usetikzlibrary{matrix,arrows,patterns,snakes,decorations.pathreplacing,3d,arrows.meta,calc}
\usepackage{etoolbox}             % Permite lógica condicional e manipulação de comandos
\usepackage{background}           % Permite adicionar conteúdo fixo no fundo da página (ex: linhas, marcas d'água)
\SetBgContents{} % Remove conteúdo padrão ("DRAFT")
\definecolor{background-color}{RGB}{233 227 206}

% --------------------------------------------------------
% GRÁFICOS E PLOTAGENS
% --------------------------------------------------------
\usepackage{pgfplots}             % Criação de gráficos vetoriais em 2D e 3D
\pgfplotsset{compat=1.18}         % Compatibilidade com versão atual do pacote
\usepackage{currfile} % Permite saber qual é o caminho completo do arquivo atual,

% --------------------------------------------------------
% LINKS E CORES
% --------------------------------------------------------
\usepackage{xcolor}              % Definição de cores (texto, links, etc.)
\usepackage[hidelinks]{hyperref} % Links sem moldura colorida no PDF
\usepackage{url}                 % Quebra automática de URLs
% Dica: se quiser links coloridos, use:
% \usepackage[colorlinks=true, linkcolor=blue, citecolor=blue, urlcolor=blue]{hyperref}

% --------------------------------------------------------
% CAIXAS E DESTAQUES VISUAIS
% --------------------------------------------------------
\usepackage[most]{tcolorbox}      % Criação de caixas coloridas e personalizadas para destaques, avisos, exemplos, etc.

% --------------------------------------------------------
% CÓDIGOS E ALGORITMOS
% --------------------------------------------------------
\usepackage{listings}            % Inclusão de código-fonte com destaque
\usepackage{algorithm}           % Ambiente de algoritmo
\usepackage[noend]{algpseudocode} % Sintaxe tipo pseudocódigo

% --------------------------------------------------------
% COMANDOS
% --------------------------------------------------------

\newcommand{\bbN}{\mathbb{N}}
\newcommand{\bbZ}{\mathbb{Z}}
\newcommand{\bbQ}{\mathbb{Q}}
\newcommand{\bbR}{\mathbb{R}}
\newcommand{\calnN}{\mathcal{N}}

\newcommand{\assunto}{ESTATÍSTICA DESCRITIVA}
\newcommand{\atividade}{lista} % prova ou lista
\newcommand{\turma}{ESTATÍSTICA}
\newcommand{\professor}{Igor Martins Silva}
\newcommand{\semestre}{2}
\newcommand{\ano}{2026}
\newcommand{\mesTexto}{outubro}
\newcommand{\mesNum}{10}
\newcommand{\dia}{14}
\newcommand{\dataTexto}{\dia\ de \mesTexto\ de \ano}
\newcommand{\dataNun}{\dia/\mesNum/\ano}

\ifdefstring{\atividade}{prova}{
    \newcommand{\titulo}{Avaliação de Estatística}
}{
    \newcommand{\titulo}{Lista de exercícios de Estatística}
}

\edef\pathCapa{\currfiledir}

\newcommand{\subtitulo}{
    %
    \noindent
    \begin{minipage}{2.8cm}
        \includegraphics[width=2.8cm]{\pathCapa Logo_Cefet.png}
    \end{minipage}
    \hfill
    \begin{minipage}{\dimexpr\textwidth-6.2cm\relax}
        \begin{center}
            CEFET-MG -- Campus Contagem \\[3pt]
            \titulo \ -- \semestre.\textordmasculine\ Semestre de \ano \\[3pt]
            \professor
        \end{center}
    \end{minipage}
    \hfill
    \begin{minipage}{3.2cm}
        \includegraphics[width=3.2cm]{\pathCapa Brasao_Brasil.png}
    \end{minipage}
    %
    
    %
    \begin{center}
        \vspace{15pt}
        \begin{tabular}{|C{250pt}|C{250pt}|}
            \hline
            \textbf{ASSUNTO} &
            \textbf{TURMA} \\
            \hline
            \assunto & \turma \\
            \hline
        \end{tabular}
    \end{center}
    %
    \vspace{15pt}
}

\newcommand{\subtituloAlt}{
    %
    \noindent
    \begin{minipage}{2.8cm}
        \includegraphics[width=2.8cm]{\pathCapa Logo_Cefet.png}
    \end{minipage}
    \hfill
    \begin{minipage}{\dimexpr\textwidth-6.2cm\relax}
        \begin{center}
            CEFET-MG -- Campus Contagem \\[3pt]
            \titulo \ -- \semestre.\textordmasculine\ Semestre de \ano \\[3pt]
            \professor
        \end{center}
    \end{minipage}
    \hfill
    \begin{minipage}{3.2cm}
        \includegraphics[width=3.2cm]{\pathCapa Brasao_Brasil.png}
    \end{minipage}
    %
    
    %
    \begin{center}
        \vspace{15pt}
        \begin{tabular}{|C{250pt}|C{100pt}|C{100pt}|}
            \hline
            \textbf{ASSUNTO} &
            \textbf{DATA} &
            \textbf{VALOR} \\
            \hline
            \assunto & \dataNun & \ValorProva\ pontos \\
            \hline
        \end{tabular}
    \end{center}
    %
    \vspace{15pt}
}

\newcommand{\valorquestao}{}

\newcommand{\definevalorquestao}{%
    \ifdefstring{\atividade}{lista}{%
        % Não faz nada para lista
    }{%
        \renewcommand{\valorquestao}{\ValorQ pontos}
    }%
}

\newenvironment{exercicioBanco}[1][]{%
    \ifdefstring{\atividade}{prova}{%
        \begin{exercicio}[#1]%
    }{%
        \begin{exercicio}%
    }%
}{%
    \end{exercicio}
}

\ifdefstring{\atividade}{lista}{
    \pagecolor{background-color}
}{
    
}



\hypersetup{
    pdfauthor={\professor},
    pdftitle={\titulo -- \assunto},
}

\title{\titulo -- \assunto}
\author{\professor}
\date{\data}

%************* INÍCIO DO DOCUMENTO *************
\begin{document}
    \subtitulo
        
    \phantomsection
    \addcontentsline{toc}{section}{Exercício 1}
    %%%%%%%%%%%%%%%%%%%%%%%%%%%%%%%%%
%
%       EXERCÍCIO
%
%%%%%%%%%%%%%%%%%%%%%%%%%%%%%%%%%

\ifdefstring{\atividade}{prova}{%
    \renewcommand{\valorquestao}{\ValorQ\ pontos}
}%

\begin{exercicioBanco}[\valorquestao]
Classifique cada uma das variáveis abaixo em qualitativa (nominal ou ordinal) ou quantitativa
(discreta ou contínua).
%
\begin{enumerate}[(a)]
    \item Ocorrência de hipertensão pré-natal em grávidas com mais de 35 anos (sim ou não).
    %
    \item Intenção de voto para presidente.
    %
    \item Perda de peso de maratonistas em uma corrida (em quilos).
    %
    \item Intensidade da perda de peso de maratonistas em uma corrida (leve, moderada, forte).
    %
    \item Grau de satisfação da população brasileira com relação ao trabalho de seu presidente (totalmente insatisfeito, insatisfeito, nem satisfeito nem insatisfeito, satisfeito, totalmente satisfeito).
\end{enumerate}
%
\end{exercicioBanco}


    \noindent\rule{\linewidth}{1pt}  % linha horizontal fina
    
    \phantomsection
    \addcontentsline{toc}{section}{Exercício 2}
    %%%%%%%%%%%%%%%%%%%%%%%%%%%%%%%%%
%
%       EXERCÍCIO
%
%%%%%%%%%%%%%%%%%%%%%%%%%%%%%%%%%

\ifdefstring{\atividade}{prova}{%
    \renewcommand{\valorquestao}{\ValorQ\ pontos}
}%

\begin{exercicioBanco}[\valorquestao]
Identifique o tipo de amostragem utilizado em cada situação:

\begin{enumerate}[(a)]
    \item Um pesquisador sorteia 30 pacientes de uma lista completa de todos os pacientes atendidos na clínica.
    \item Em um hospital, seleciona-se um paciente a cada 10 registros da agenda.
    \item Os pacientes são separados por faixa etária e, em seguida, são sorteados pacientes dentro de cada faixa.
    \item São sorteadas algumas clínicas de fisioterapia e todos os pacientes dessas clínicas são avaliados.
\end{enumerate}
\end{exercicioBanco}

    \noindent\rule{\linewidth}{1pt}  % linha horizontal fina
    
    \phantomsection
    \addcontentsline{toc}{section}{Exercício 3}
    %%%%%%%%%%%%%%%%%%%%%%%%%%%%%%%%%
%
%       EXERCÍCIO
%
%%%%%%%%%%%%%%%%%%%%%%%%%%%%%%%%%

\ifdefstring{\atividade}{prova}{%
    \renewcommand{\valorquestao}{\ValorQ\ pontos}
}%

\begin{exercicioBanco}[\valorquestao]
Um estudo registrou o tempo de permanência (em dias) de $10$ pacientes internados em uma ala hospitalar:
\[ 5, \; 7, \; 3, \; 12, \; 7, \; 8, \; 15, \; 7, \; 10, \; 6 \]

Com base nesses dados, determine:

\begin{enumerate}[(a)]
    \item A média, a mediana e a moda do tempo de internação.
    \item O primeiro quartil ($Q_1$) e o terceiro quartil ($Q_3$).
    \item Se o valor discrepante $15$ fosse substituído por $25$, qual das três medidas do item (a) sofreria maior alteração? Justifique brevemente.
\end{enumerate}
\end{exercicioBanco}

    \noindent\rule{\linewidth}{1pt}  % linha horizontal fina
    
    \phantomsection
    \addcontentsline{toc}{section}{Exercício 4}
    %%%%%%%%%%%%%%%%%%%%%%%%%%%%%%%%%
%
%       EXERCÍCIO
%
%%%%%%%%%%%%%%%%%%%%%%%%%%%%%%%%%

\definevalorquestao

\begin{exercicioBanco}[\valorquestao]
A tabela a seguir apresenta a distribuição de frequências relativas aos tempos de atendimento (em minutos) registrados em um posto de saúde:

\begin{center}
\begin{tabular}{|c|c|}
\hline
Tempo (min) & Frequência ($f_i$) \\ \hline
$0 \vdash 10$ & 3 \\ \hline
$10 \vdash 20$ & 5 \\ \hline
$20 \vdash 30$ & 8 \\ \hline
$30 \vdash 40$ & 4 \\ \hline
$40 \vdash 50$ & 2 \\ \hline
\end{tabular}
\end{center}

Com base nessa tabela, classifique os dados quanto sua assimetria, usando o 2.\textordmasculine\ coeficiente de assimetria de Pearson, \(e_2\). Para isso, considere a seguinte classificação:
\begin{center}
    \renewcommand{\arraystretch}{1.2}
    \begin{tabular}{|c|c|}
        \hline
        \textbf{Valor de } $e_2$ & \textbf{Classificação da Assimetria} \\ \hline
        $|e_2| = 0$ & Simétrica \\ \hline
        $0 < |e_2| \le 0{,}15$ & Assimetria fraca \\ \hline
        $0{,}15 < |e_2| \le 1$ & Assimetria moderada \\ \hline
        $|e_2| > 1$ & Assimetria forte \\ \hline
    \end{tabular}
\end{center}
\end{exercicioBanco}

    \noindent\rule{\linewidth}{1pt}  % linha horizontal fina
    
    \phantomsection
    \addcontentsline{toc}{section}{Exercício 5}
    %%%%%%%%%%%%%%%%%%%%%%%%%%%%%%%%%
%
%       EXERCÍCIO
%
%%%%%%%%%%%%%%%%%%%%%%%%%%%%%%%%%

\ifdefstring{\atividade}{prova}{%
    \renewcommand{\valorquestao}{\ValorQ\ pontos}
}%

\begin{exercicioBanco}[\valorquestao]
A tabela a seguir apresenta a distribuição de frequências dos salários (em salários mínimos) dos funcionários de uma empresa:

\begin{center}
\begin{tabular}{|c|c|}
\hline
Salário ($S_i$) & Frequência ($f_i$) \\ \hline
$2 \vdash 4$ & 4 \\ \hline
$4 \vdash 6$ & 10 \\ \hline
$6 \vdash 8$ & 4 \\ \hline
$8 \vdash 10$ & 2 \\ \hline
\end{tabular}
\end{center}

Com base nesses dados, determine:

\begin{enumerate}[(a)]
    \item A variância populacional ($\sigma^2$).
    \item O desvio padrão populacional ($\sigma$).
    \item O coeficiente de variação ($CV$).
\end{enumerate}
\end{exercicioBanco}

    \noindent\rule{\linewidth}{1pt}  % linha horizontal fina
    
    \phantomsection
    \addcontentsline{toc}{section}{Exercício 6}
    %%%%%%%%%%%%%%%%%%%%%%%%%%%%%%%%%
%
%       EXERCÍCIO
%
%%%%%%%%%%%%%%%%%%%%%%%%%%%%%%%%%

\ifdefstring{\atividade}{prova}{%
    \renewcommand{\valorquestao}{\ValorQ\ pontos}
}%

\begin{exercicioBanco}[\valorquestao]
Um levantamento estatístico, feito em determinada região do país, mostrou que jovens com idades entre 4 e 17 anos assistem à televisão, em média, durante 6 horas por dia ($\mu = 6$). A tabela a seguir apresenta outras estatísticas produzidas por esse levantamento:

\begin{center}
\begin{tabular}{|c|c|}
\hline
\textbf{Estatística} & \textbf{Tempo ($T$, em horas)} \\ \hline
1.º quartil ($Q_1$) & 2 \\ \hline
2.º quartil ($Q_2$) & 4 \\ \hline
3.º quartil ($Q_3$) & 8 \\ \hline
1.º decil ($D_1$) & 1 \\ \hline
9.º decil ($D_9$) & 10 \\ \hline
\end{tabular}
\end{center}

Com base estritamente nas informações fornecidas na tabela, determine:

\begin{enumerate}[(a)]
    \item O coeficiente quartílico de assimetria de Bowley e classifique a distribuição quanto à assimetria.
    \item O coeficiente de curtose baseado em quartis e decis e classifique a curva de frequência quanto à curtose.
\end{enumerate}
\end{exercicioBanco}

    \noindent\rule{\linewidth}{1pt}  % linha horizontal fina
    
    \phantomsection
    \addcontentsline{toc}{section}{Exercício 7}
    %%%%%%%%%%%%%%%%%%%%%%%%%%%%%%%%%
%
%       EXERCÍCIO
%
%%%%%%%%%%%%%%%%%%%%%%%%%%%%%%%%%

\definevalorquestao

\begin{exercicioBanco}[\valorquestao]

O histograma representa a distribuição das estaturas de 100 pessoas e as respectivas frequências. Por exemplo, na 3ª classe (\(155\text{--}160\)) estão situadas \(11\%\) das pessoas com estatura de \(1,55\text{ m}\) a \(1,59\text{ m}\). A 5ª classe (\(165\text{--}170\)) chama-se classe mediana. Pelo ponto \(M\) situado na classe mediana, traça-se uma reta paralela ao eixo das frequências, de modo a dividir a área da figura formada pelos nove retângulos das frequências em duas regiões de mesma área. Determine a abscissa do ponto \(M\) (mediana das observações).

\begin{center}
\begin{tikzpicture}
\definecolor{azulAba}{RGB}{133,167,220}

% Eixo Principal (Esquerdo)
\begin{axis}[
    x=0.22cm,
    height=6.5cm,
    ymin=0, ymax=28,
    xmin=145, xmax=190,
    axis lines=left,
    clip=false,
    ytick={6,8,11,15,24},
    xtick={145,150,155,160,165,170,175,180,185,190},
    ylabel={frequência (\%)},
    every axis y label/.style={at={(axis description cs:0,1)},anchor=south west,font=\small},
    y tick label style={font=\scriptsize},
    x tick label style={font=\scriptsize, anchor=north},
    % Controle rigoroso de camadas e linhas horizontais
    set layers,
    ymajorgrids=true,
    grid style={gray!50, thin},
    ybar,
    bar width=1.1cm,
    bar shift=0.55cm,
    enlargelimits=false
]

% Barras azuis opacas (ficam por cima e cobrem as linhas atrás delas)
\addplot[fill=azulAba, draw=white, thick] coordinates {
    (145, 6)
    (150, 8)
    (155, 11)
    (160, 15)
    (165, 24)
    (170, 14)
    (175, 10)
    (180, 7)
    (185, 5)
};

% Linha tracejada da Mediana M e rótulo M
\draw[dashed, black, thick] (axis cs:167.08, 0) -- (axis cs:167.08, 27);
\node[below, font=\small] at (axis cs:167.08, -1.8) {M};

\end{axis}

% Eixo Y Secundário (Direito)
\begin{axis}[
    x=0.22cm,
    height=6.5cm,
    ymin=0, ymax=28,
    xmin=145, xmax=190,
    axis y line=right,
    axis x line=none,
    ytick={5,7,10,14},
    ylabel={frequência (\%)},
    every axis y label/.style={at={(axis description cs:1,1)},anchor=south east,font=\small},
    y tick label style={font=\scriptsize},
    % Linhas horizontais do lado direito
    set layers,
    ymajorgrids=true,
    grid style={gray!50, thin},
    enlargelimits=false
]
\end{axis}
\end{tikzpicture}
\end{center}

\end{exercicioBanco}

    \noindent\rule{\linewidth}{1pt}  % linha horizontal fina
    
    \phantomsection
    \addcontentsline{toc}{section}{Exercício 8}
    %%%%%%%%%%%%%%%%%%%%%%%%%%%%%%%%%
%
%       EXERCÍCIO
%
%%%%%%%%%%%%%%%%%%%%%%%%%%%%%%%%%

\ifdefstring{\atividade}{prova}{%
    \renewcommand{\valorquestao}{\ValorQ\ pontos}
}%

\begin{exercicioBanco}[\valorquestao]
O que acontece com a mediana, a média e o desvio-padrão de um conjunto de dados quando:
%
\begin{enumerate}[(a)]
    \item o valor 10 é somado a cada observação?
    %
    \item cada observação é multiplicada por 2?
    %
    \item subtrai-se a média de cada observação?
    %
    \item de cada observação, subtrai-se a média e divide-se pelo desvio-padrão?
\end{enumerate}
%
\end{exercicioBanco}


    \noindent\rule{\linewidth}{1pt}  % linha horizontal fina
    
    \phantomsection
    \addcontentsline{toc}{section}{Exercício 9}
    %%%%%%%%%%%%%%%%%%%%%%%%%%%%%%%%%
%
%       EXERCÍCIO
%
%%%%%%%%%%%%%%%%%%%%%%%%%%%%%%%%%

\ifdefstring{\atividade}{prova}{%
    \renewcommand{\valorquestao}{\ValorQ\ pontos}
}%

\begin{exercicioBanco}[\valorquestao]
O volume diário de leite produzido por uma cooperativa agrícola (em milhares de litros) foi registrado durante 40 dias consecutivos e é apresentado no rol a seguir:

\begin{center}
    \renewcommand{\arraystretch}{1.3}
    \begin{tabular}{|C{30pt}|C{30pt}|C{30pt}|C{30pt}|C{30pt}|C{30pt}|C{30pt}|C{30pt}|C{30pt}|C{30pt}|}
        \hline
        10.2 & 10.4 & 10.5 & 10.7 & 11.0 & 11.1 & 11.3 & 11.4 & 11.5 & 11.5 \\
        \hline
        11.6 & 11.6 & 11.7 & 11.7 & 11.8 & 11.8 & 11.8 & 11.9 & 11.9 & 12.0 \\
        \hline
        12.0 & 12.1 & 12.1 & 12.2 & 12.2 & 12.3 & 12.3 & 12.4 & 12.4 & 12.5 \\
        \hline
        12.6 & 12.7 & 12.8 & 13.0 & 13.1 & 13.3 & 13.5 & 13.6 & 13.8 & 14.2 \\
        \hline
    \end{tabular}
\end{center}

\begin{enumerate}[(a)]
    \item Construa o histograma dessa variável estatística, considerando intervalos de classe com amplitude igual a \(0.8\), iniciando a primeira classe em \(10.0\).
    %
    \item Construa o boxplot (diagrama de caixas) correspondente a esse conjunto de dados.
\end{enumerate}
%
\end{exercicioBanco}


    \noindent\rule{\linewidth}{1pt}  % linha horizontal fina
    
    \phantomsection
    \addcontentsline{toc}{section}{Exercício 10}
    %%%%%%%%%%%%%%%%%%%%%%%%%%%%%%%%%
%
%       EXERCÍCIO
%
%%%%%%%%%%%%%%%%%%%%%%%%%%%%%%%%%

\definevalorquestao

\begin{exercicioBanco}[\valorquestao]
O quadro a seguir apresenta a quantidade de um tipo de pão vendido em uma semana em uma padaria.

\begin{longtable}{l c}
\hline
Dia da semana & Número de pães vendidos \\
\hline
Domingo       & 250 \\
Segunda-feira & 208 \\
Terça-feira   & 215 \\
Quarta-feira  & 251 \\
Quinta-feira  & 187 \\
Sexta-feira   & 187 \\
Sábado        & 186 \\
\hline
\end{longtable}

O dono da padaria decidiu que, na semana seguinte, a produção diária desse tipo de pão seria igual ao número de pães vendidos no dia da semana em que tal quantidade foi a mais próxima da média das quantidades vendidas na semana.

Qual o dia da semana utilizado como referência para a quantidade de pães a serem produzidos diariamente?
\end{exercicioBanco}


    \noindent\rule{\linewidth}{1pt}  % linha horizontal fina
    
    \phantomsection
    \addcontentsline{toc}{section}{Exercício 11}
    %%%%%%%%%%%%%%%%%%%%%%%%%%%%%%%%%
%
%       EXERCÍCIO
%
%%%%%%%%%%%%%%%%%%%%%%%%%%%%%%%%%

\ifdefstring{\atividade}{prova}{%
    \renewcommand{\valorquestao}{\ValorQ\ pontos}
}%

\begin{exercicioBanco}[\valorquestao]
O gráfico apresenta a taxa de desemprego (em \%) para o período de março de 2008 a abril de 2009, com base nos dados observados nas regiões metropolitanas de Recife, Salvador, Belo Horizonte, Rio de Janeiro, São Paulo e Porto Alegre.

\begin{center}
\begin{tikzpicture}
\begin{axis}[
    x=0.85cm, % Estica a distância entre cada ponto no eixo X (ajuste esse valor se quiser mais/menos esticado)
    height=6cm, % Define a altura do eixo Y
    xlabel={Meses (Mar/08 a Abr/09)},
    ylabel={Taxa de desemprego (\%)},
    ymin=6.5, ymax=9.8,
    xtick={1,...,14},
    xticklabels={
        Mar/08,Abr/08,Mai/08,Jun/08,Jul/08,Ago/08,Set/08,Out/08,
        Nov/08,Dez/08,Jan/09,Fev/09,Mar/09,Abr/09},
    x tick label style={rotate=45, anchor=east, font=\scriptsize},
    y tick label style={font=\scriptsize},
    label style={font=\small},
    grid=both,
    major grid style={gray!30},
    minor grid style={gray!15}
]
\addplot[
    color=blue,
    mark=*,
    thick,
    nodes near coords,
    point meta=explicit symbolic,
    every node near coord/.append style={font=\scriptsize, anchor=south}
] coordinates {
    (1,8.6) [8,6]
    (2,8.5) [8,5]
    (3,7.9) [7,9]
    (4,7.9) [7,9]
    (5,8.1) [8,1]
    (6,7.6) [7,6]
    (7,7.7) [7,7]
    (8,7.5) [7,5]
    (9,7.6) [7,6]
    (10,6.8) [6,8]
    (11,8.2) [8,2]
    (12,8.5) [8,5]
    (13,9.0) [9,0]
    (14,8.9) [8,9]
};
\end{axis}
\end{tikzpicture}
\end{center}

A mediana dessa taxa de desemprego, no período de março de 2008 a abril de 2009, foi de quanto?

\end{exercicioBanco}


    \noindent\rule{\linewidth}{1pt}  % linha horizontal fina
    
    \phantomsection
    \addcontentsline{toc}{section}{Exercício 12}
    %%%%%%%%%%%%%%%%%%%%%%%%%%%%%%%%%
%
%       EXERCÍCIO
%
%%%%%%%%%%%%%%%%%%%%%%%%%%%%%%%%%

\definevalorquestao

\begin{exercicioBanco}[\valorquestao]
Um posto de saúde registrou a quantidade de vacinas aplicadas contra febre amarela nos últimos cinco meses:
\begin{itemize}
    \begin{multicols}{3}
        \item 1º mês: 21;
        \item 2º mês: 22;
        \item 3º mês: 25;
        \item 4º mês: 31;
        \item 5º mês: 21.
    \end{multicols}
\end{itemize}

No início do primeiro mês, esse posto de saúde tinha 228 vacinas em estoque. A política de reposição prevê a aquisição de novas vacinas, no início do sexto mês, de tal forma que a quantidade inicial em estoque para os próximos meses seja igual a 12 vezes a média das quantidades mensais aplicadas nos últimos cinco meses.

Qual a quantidade de vacinas que deve ser adquirida no início do sexto mês?
\end{exercicioBanco}


    \noindent\rule{\linewidth}{1pt}  % linha horizontal fina
    
\end{document}
%*************** FINAL DO DOCUMENTO ***************
